\bookchapter{附錄C：極簡術語「黑話」對照表}

\begin{center}
\begin{tabularx}{\linewidth}{|X|X|X|}
\hline
\textbf{你可能聽到嘅詞} & \textbf{其實就係指……} & \textbf{一句話解釋} \tabularnewline
\hline
SoC / 晶片 & 手機嘅大腦同心臟 & 決定手機快唔快、遊戲卡唔卡嘅核心部件。 \tabularnewline
RAM (記憶體) & 手機嘅「短期工作桌」 & 桌越大，可以同時開嘅App就越多，切換越流暢。 \tabularnewline
ROM (存儲) & 手機嘅「大倉庫」 & 存相、影片、App嘅地方。128G、256G講嘅就係佢。 \tabularnewline
OLED / LCD & 螢幕嘅兩種材質 & OLED色彩更艷、更省電；LCD更護眼、成本低。 \tabularnewline
NFC & 近距離無線通訊 & 等你手機可以變身公交卡、門禁卡、銀行卡嘅「貼貼」功能。 \tabularnewline
OTA更新 & 無線系統升級 & 喺【設定】入面撳嘅嗰個「系統更新」，就係佢。 \tabularnewline
Root & 安卓超級管理員權限 & 獲得後可以修改系統檔案，但都有風險。 \tabularnewline
BL鎖 & Bootloader鎖 & 廠商加嘅安全鎖，解鎖後先可以刷機。 \tabularnewline
NPU & 神經網絡處理器 & 專門處理AI任務嘅晶片，比CPU快幾十倍、功耗更低。 \tabularnewline
端側AI & 手機本地運行嘅AI & AI喺手機上直接運行，唔需要聯網，速度快、私隱好。 \tabularnewline
Qi2 & 新一代無線充電標準 & 帶磁吸定位嘅無線充電，類似MagSafe，充電更方便。 \tabularnewline
eSIM & 嵌入式SIM卡 & 冇實體卡嘅SIM，直接寫入手機，換運營商唔使換卡。 \tabularnewline
LTPO & 自適應刷新率技術 & 螢幕刷新率根據內容自動調節（1-120Hz），流暢又省電。 \tabularnewline
\hline
\end{tabularx}
\end{center}
